\section{Current Situation}

\subsection{Problem Statement}

According to UNICEF, between the years 2005 and 2022, over 105,000 children have been recruited to participate in armed conflicts. While it is widely known that these children are oftentimes used as fighters, less people are aware of the fact that they are occasionally also used as cooks, spies, sex slaves, bomb makers, guards, porters, messengers, and suicide bombers\footnote{UNICEF. \emph{Children Recruited by Armed Forces or Armed Groups}. 2021. \url{https://www.unicef.org/protection/children-recruited-by-armed-forces}}. At the moment, this practice is particularly prevalent across certain parts of Asia, Africa, and South America\footnote{EBSCO Information Services. \emph{Child Soldiers}. Research Starters: Military History and Science. \url{https://www.ebsco.com/research-starters/military-history-and-science/child-soldiers\#full-article}}. More importantly, Human Rights Watch affirms that around 30\% of armed groups integrated by children contain girls, generally subject to sexual abuse and sometimes even taken as mistresses of rebel army leaders in countries like Uganda, Angola or Sierra Leone\footnote{Human Rights Watch. \emph{Stop the Use of Child Soldiers - Statements and Resolutions}. 2001. \url{https://www.hrw.org/news/2001/06/01/stop-use-child-soliders-statements-and-resolutions}}. There have also been cases of child soldier recruitment in the Russia-Ukraine and Israel-Palestine Conflicts. For example, according to The Child Labour Coalition, in certain parts of Siberia, minors are put in charge of sewing items and preparing care packages for Russian troops fighting in Ukraine and city administrators in Krasnoyarsk have sent representatives to some schools to recruit children\footnote{Greenberger Sydney. \emph{New Report on Children and Armed Conflict — The Use of Child Soldiers}. The Child Soldier Coalition. 2024. \url{https://stopchildlabor.org/new-report-on-children-and-armed-conflict-the-use-of-child-soldiers/}}. Similarly, a very recent Report of the Secretary-General on Children and Armed Conflict drafted this year has verified that Palestinian children were used by Israeli armed and security forces as human shields during operations in the West Bank as well as in the Gaza Strip\footnote{United Nations General Assembly. \emph{Children and armed conflict - Report of the Secretary-General (A/80/723-S/2026/357)}. 2026. \url{https://reliefweb.int/report/world/children-and-armed-conflict-report-secretary-general-a80723-s2026357}}.

A lot of children are violently taken away from their homes to serve as soldiers or to carry out other roles in armed conflicts\footnote{International Committee of the Red Cross. \emph{Children Associated with Armed Forces or Armed Groups}. Geneva, Switzerland. 2013. \url{https://www.icrc.org/sites/default/files/external/doc/en/assets/files/publications/icrc-002-0824.pdf}}. However, many of them also decide to join armed groups because they are indoctrinated through nationalist arguments or because they see this as an opportunity to escape poverty. In fact, a report developed by War Child UK titled “Tug of War: Children in Armed Groups in DRC” exposes how a large part of children in the Democratic Republic of Congo voluntarily enlist in the military. According to this report, what makes them willing to become involved in conflict is the absence of a stable and caring family environment\footnote{\emph{Why do Children Become Child Soldiers?} War Child. 2018. \url{https://www.warchild.org.uk/news/why-do-children-become-child-soldiers}}. This suggests that there is a pressing need to tackle poverty and instability across the regions where the problem of child soldiers is still present.

In addition, it is also worth considering how, when taking part in armed conflicts, children are often forced to witness very violent killings, rapes, beheadings, limb amputations, and burnings of people alive. Along these lines, many children who are involved in armed conflicts from a young age are also indoctrinated by the armed groups that recruit them. All of this leaves them with lasting physical and mental health problems which might affect their capacity to reintegrate themselves back into society once they are released\footnote{Kaplan Eben. \emph{Child Soldiers around the World}. Council on Foreign Relations. 2005. \url{https://www.cfr.org/backgrounders/child-soldiers-around-world}}. Furthermore, child soldiers are also deprived of a quality education and denied the possibility of living a normal childhood. This limits their future professional opportunities and perpetuates cycles of poverty in post-conflict regions. It is therefore necessary to further expand and improve existing rehabilitation and reintegration schemes, tailoring them to the circumstances of each former child soldier, focusing, in particular, on the activities that they had to perform while recruited by armed groups during periods of war.

Another issue which has generally been overlooked is the fact that girls often find it harder to escape and usually encounter greater social stigma once they are able to return home\footnote{Plan International. \emph{Hidden in plain sight: girls associated with armed groups}. Plan International. 2026. \url{https://plan-international.org/blog/2026/02/11/hidden-in-plain-sight-girls-associated-with-armed-groups/}}. Researcher Milfrid Tonheim exposes how most of the child soldiers that are freed are boys. She explains how this is because international organizations such as UNICEF need to cooperate with the commanders of armed groups in order to free children who have been forcibly recruited. They usually receive lists with the names of the underaged children which they need to save but the girl’s names are most of the time not included since the commanders of the armed groups that take them generally see them as their mistresses. On top of this, many of the girls who manage to escape are then marginalised by their communities. Some families find it hard to accept that their daughters who were once kidnapped and raped by commanders of rebel armed forces are no longer virgins and might not be able to get married. If the daughters arrive pregnant things get even worse, since they and their newborns will be considered financial burdens\footnote{Grønhaug Kristene. \emph{The Number of Child Soldiers in the World is Increasing – Almost Half of Them are Girls}. Flykninghjelpen Norwegian Refugee Council. \url{https://www.nrc.no/shorthand/stories/child-soldiers/index.html}}.

Finally, although a number of treaties, conventions, and resolutions have been put in place in order to tackle the use and recruitment of child soldiers in conflict zones, these are rarely ratified by all countries and are oftentimes difficult to enforce. Armed groups are aware that children can be very easily lured into becoming involved in the military and they effectively take advantage of this. Hence, it is also important to strengthen the enforcement mechanisms that are already in place, while making sure that both governments and non-state groups who use and recruit child soldiers are adequately sanctioned and held accountable for their actions.

\subsection{Possible Solutions}

This section examines how already existing frameworks and initiatives could be improved in order to better tackle the problem of child soldier use and recruitment across the globe. It also puts forward some possible solutions which have not yet been implemented and which we would like all delegates to expand further in their draft resolutions.

Solutions to the child soldier problem cannot sit in one phase of the conflict cycle. They have to run from prevention, through release and reintegration, to accountability, and they have to work while fighting is still active, not only after a peace agreement. Delegates should keep proposals within DISEC’s disarmament and security mandate rather than straying into field protection or prosecution, which sit with other organs.

\subsubsection{Prevention and root causes}

Recruitment is driven by poverty and the absence of alternatives as much as by force. War Child UK’s research in the Kivus found that many children enlist ‘’voluntarily’’ because they lack a stable family environment and see no economic future\footnote{War Child UK. (2018). \emph{Why do children become child soldiers?}\url{https://www.warchild.org.uk/news/why-do-children-become-child-soldiers}}. Prevention therefore means restoring access to education, one of the strongest protective factors\footnote{Education Above All. (n.d.). \emph{Education of former child soldiers.}\url{https://www.educationaboveall.org/sites/default/files/research/attachments/02.\%20EducationChildSoldiersEN.pdf}}, expanding adolescent livelihoods, and strengthening birth and civil registration so a child's age can actually be verified at recruitment, now a standard commitment in UN action plans\footnote{Office of the Special Representative of the Secretary-General for Children and Armed Conflict. (n.d.). \emph{Action plans.}\url{https://childrenandarmedconflict.un.org/tools-for-action/action-plans/}}. The Safe Schools Declaration and Security Council Resolution 2601 (2021) reduce exposure by limiting the military use of schools\footnote{Security Council Report. (2026, June). \emph{Children and armed conflict: June 2026 monthly forecast.}\url{https://www.securitycouncilreport.org/monthly-forecast/2026-06/children-and-armed-conflict-12.php}}.

\subsubsection{Adapting DDR}

The IDDRS already treats child release as a human rights obligation that is not contingent on any political negotiation and should be pursued during conflict\footnote{United Nations Inter-Agency Working Group on DDR. (2023). \emph{IDDRS 5.20: Children and DDR.}\url{https://www.unddr.org/wp-content/uploads/2023/03/IDDRS-5.20-Children-and-DDR.pdf}}, and the 2019 UN Approach to DDR extends guidance to settings with no peace agreement\footnote{United Nations. (2019). \emph{The IDDRS and the UN approach to DDR.}\url{https://www.unddr.org/the-iddrs/}}. Two fixes matter most: decouple DDR eligibility from surrendering a weapon, since that criterion excludes girls and support roles\footnote{Global Coalition for Reintegration of Child Soldiers. (2020). \emph{Reframing child reintegration: From humanitarian action to development.}\url{https://childrenandarmedconflict.un.org/wp-content/uploads/2020/09/GCR-Reframing-Child-Reintegration-92020.pdf}}, and fund community-based reintegration over several years rather than in short humanitarian cycles. Funding fell by roughly a third between 2012 and 2016 as need rose\footnote{Office of the Special Representative of the Secretary-General for Children and Armed Conflict. (n.d.). \emph{Global Coalition for Reintegration of Child Soldiers.}\url{https://childrenandarmedconflict.un.org/en/global-coalition-for-reintegration-of-child-soldiers}}, so the Global Coalition for Reintegration argues for treating reintegration as development, not short-term relief.

\subsubsection{Gender}

Girls are undercounted because releases run through commander-supplied lists that leave them off, so those freed are disproportionately boys\footnote{Grønhaug, K. (n.d.). \emph{The number of child soldiers in the world is increasing – almost half of them are girls.} Norwegian Refugee Council.\url{https://www.nrc.no/shorthand/stories/child-soldiers/index.html}}, and returnees face the worst stigma, girls' mothers most of all. Resolution 2427 (2018) already calls for tailored programmes\footnote{United Nations Security Council. (2018). \emph{Resolution 2427.}\url{https://unscr.com/en/resolutions/doc/2427/}}, making that real means counting girls in negotiated releases, funding gender specific psychosocial and medical care, and investing in community sensitisation.

\subsubsection{Accountability}

The gap is enforcement, not law. Strengthen the Monitoring and Reporting Mechanism and the action plan system, of which 38 have been signed, and 12 fully complied with, and tie delisting to verified reintegration, rather than a declared halt\footnote{Office of the Special Representative of the Secretary-General for Children and Armed Conflict. (n.d.). \emph{Action plans.}\url{https://childrenandarmedconflict.un.org/tools-for-action/action-plans/}}. Despite this architecture, the Secretary-General verified 41,370 grave violations in 2024, up 25\%\footnote{United Nations Secretary-General. (2025). \emph{Children and armed conflict: Report of the Secretary-General (covering 2024).}\url{https://childrenandarmedconflict.un.org/wp-content/uploads/2025/06/Secretary-General-Annual-Report-on-Children-and-Armed-Conflict-Covering-2024.pdf}}. The ICC has real precedent under Article 8 of the Rome Statute in Lubanga, Ntaganda and Ongwen.\footnote{International Criminal Court. \emph{The Prosecutor v. Thomas Lubanga Dyilo} (14 years, 2012); \emph{Ntaganda} (30 years, 2019); \emph{Ongwen} (25 years, 2021).\url{https://www.icc-cpi.int/drc/lubanga}} Since leaders evade any single court, states should use universal jurisdiction and national prosecutions, while distinguishing child victims from adult perpetrators, the dilemma the Ongwen case exposed. National leverage like the US Child Soldiers Prevention Act works only when applied consistently rather than waived for strategic partners, as noted in Key Stakeholders.
